\documentclass[11pt]{article}
\usepackage[T1]{fontenc}
\usepackage[utf8]{inputenc}
\usepackage{lmodern,microtype,geometry,amsmath,amssymb,amsthm,mathtools}
\usepackage[colorlinks=true,linkcolor=blue,citecolor=blue,urlcolor=blue]{hyperref}
\geometry{margin=1in}
\setlength{\emergencystretch}{2em}
\hypersetup{pdftitle={Online learning of AND under Gaussian neural-network priors},pdfauthor={ChatGPT-6 Astra; Claude; Kaarel Hanni; Dmitry Vaintrob}}
\newtheorem{theorem}{Theorem}
\newtheorem{lemma}[theorem]{Lemma}
\newtheorem{proposition}[theorem]{Proposition}
\newtheorem{corollary}[theorem]{Corollary}
\newtheorem{remark}[theorem]{Remark}
\title{Online learning of AND under Gaussian neural-network priors}
\author{ChatGPT-6 Astra \and Claude (Anthropic) \and Kaarel H\"anni \and Dmitry Vaintrob}
\date{Research draft, 30 October 2026}
\begin{document}
\maketitle
\begin{abstract}
The conjunction of all input bits can force exponentially many online
mistakes for a deterministic learner that selects a smallest consistent
Boolean circuit. Does the same target defeat exact Bayesian prediction
with Gaussian neural-network priors? We prove two contrasting results
for fully connected tanh networks at the weak-prior scaling $\gamma=0$,
retaining a global output-margin restriction and exact label conditioning.
At each fixed depth, the conjunction has a polynomial negative log
function probability, so posterior sampling makes only polynomially many
expected mistakes on every input order. Posterior majority and a fixed
confidence criterion obey corresponding bounds. These statements hold
in the Gaussian-process width limit and at every sufficiently large
finite width. In contrast, at depth $O(d)$ the function prior becomes
nearly uniform on all Boolean functions: at every sufficiently large
width, every ordering of the $2^d$ inputs incurs at least
$(2/5)2^d$ expected sampled mistakes, and every new input fails the
strict $2/3$ confidence criterion. This lower bound does not establish
majority-prediction errors. An arbitrarily small first-layer variance anisotropy gives a separate
strong-prior exponential lower bound, for AND and for a unit-weight
one-neuron majority target. Breaking exact variance symmetry is essential to
that construction. We also analyze why the minimum-circuit conjunction
argument does not identify the minimum-energy neural functions. No counterexample for the equal-variance strong-prior or
$\gamma\ge1$ regimes is proved here.
\end{abstract}

\section{The learning question and the model}
The target is the Boolean conjunction
$f_\wedge(x)=+1$ only at $x=(1,\ldots,1)$, on
$\mathcal X=\{-1,1\}^d$; it is negative elsewhere. Its minimum
binary AND/OR/NOT circuit size is $d-1$. An online learner sees an
input, predicts a label, and receives its true label. Every observation
is retained. Inputs may be chosen using the entire past, but before
the current private random prediction. A fixed input permutation is
an important special case.

For deterministic minimum-circuit selection, a conjunction ladder can
force $2^d-d-2$ mistakes~\cite{occam}. After all inputs of Hamming weight
below $k$ have been labeled negative, and the all-ones input positive,
the smallest consistent circuits are $k$-variable conjunctions. Each
has a different remaining weight-$k$ positive input. A query at the
selected conjunction's input therefore causes a mistake. Fresh uniform
randomization over those circuits changes that conclusion. A neural
posterior, however, weights functions by parameter-space probability,
not gate count, so neither the deterministic construction nor its
uniform-tie remedy transfers without a new argument.

We use $L-1$ hidden layers of common width $N$, tanh activations, a
linear readout, and one appended constant input. There are no biases
in later layers. The independent centered Gaussian weight variances
in the one-parameter family are
\begin{equation}\label{eq:model-gamma}
 \operatorname{Var}(W^1_{ij})=N^{1-\gamma}/d,\qquad
 \operatorname{Var}(W^s_{ij})=N^{-\gamma}\ (1<s<L),\qquad
 \operatorname{Var}(v_i)=N^{-\gamma-1}.
\end{equation}
The main positive and negative theorems of this note concern
$\gamma=0$. Before observing labels, condition on the label-independent
event $|z_w(x)|>a$ for every input; the default is $a=1/2$.
For a function $g$, write
\[
 U_g=\{w:g(x)z_w(x)>a\ \forall x\},\qquad
 \pi_N(g)=\frac{\mathbb P(U_g)}{\mathbb P(\bigcup_hU_h)}.
\]
Training conditions this fixed function prior on all observed labels.
Nothing flattens function ranks, changes mass ratios, or replaces the
posterior with an optimizer. A separately labeled final subsection
changes the first-layer variances by arbitrarily small unequal factors
and proves a strong-prior lower bound for that perturbed prior.

When a bounded-weight teacher complexity is used below, $C_L(f)$
denotes the least hidden width of a depth-$L$ tanh implementation of
$f$ whose coefficients all have magnitude at most that width and whose
signed outputs exceed $1/2$ everywhere. Thus $N\ge C_L(f)$ is a
specific sufficient-width convention, not merely representability by
unbounded real weights. The anisotropic majority example satisfies it
with $C_2(f)=N=1$.

A sampled mistake means that a fresh posterior function gives the
wrong label. A majority mistake means that the more probable label is
wrong, with any stated tie rule. A confidence failure means that the
correct label has probability at most $2/3$. Nearly fair predictions
can cause many sampling and confidence errors while their slight
majority is correct; we keep these claims separate.

The order of limits is explicit throughout: first fix $d$ and a finite
$L$, then let $N$ grow. Choosing $L=L_d=O(d)$ in the lower bound is
compatible with this order. The finite-width thresholds are not
quantified. In particular, none of these width-limit statements
interchanges the strong-prior limit, which instead makes Gaussian
variance tend to zero at each fixed width.

\section{A polynomial bound for AND at fixed depth}
\label{gpupper:section}

The wide-network Gaussian limit also gives a positive result that
restricts where an AND lower bound can occur. Fix $d\ge1$ and $L\ge2$,
write $m=2^d$ and $r=L-1$, and let
\[
 \begin{gathered}
 f_\wedge(x)=+1\ \Longleftrightarrow\ x=(1,\ldots,1),
 \qquad \Psi(t)=\frac2\pi\arcsin t,\\
 K_r(x,y)=\Psi^{\circ r}\!\left(\frac{1+x\cdot y}{d+1}\right).
 \end{gathered}
\]
Consider $G\sim\mathcal N(0,K_r)$, conditioned before observing labels
on $|G(x)|>a$ at every Boolean input. We take $a=1/2$ in the neural
application. This is the $\gamma=0$ limiting function prior of the actual
layer-scaled Gaussian tanh network; the nonpositive-exponent
posterior-limit proposition in~\cite{nn-subcritical} supplies this limit.

\begin{theorem}[Fixed-depth AND target mass]\label{gpupper:mass}
Put
\begin{equation}\label{gpupper:budget}
 A=a+\sqrt{2\log(4m)},\qquad
 D_{d,L,a}=A^2(d+1)(d^2-d+1)\left(\frac\pi2\right)^r.
\end{equation}
The unconditioned Gaussian probability of the target's full margin
event satisfies
\begin{equation}\label{gpupper:rawmass}
 \mathbb P\{f_\wedge(x)G(x)>a\text{ for every }x\}
 \ge\frac14 e^{-D_{d,L,a}}.
\end{equation}
Consequently the globally margin-conditioned function prior has
\begin{equation}\label{gpupper:targetcost}
 -\log\pi_\infty(f_\wedge)\le D_{d,L,a}+\log4.
\end{equation}
At fixed $L$ and fixed $a$, the right side is polynomial in $d$.
\end{theorem}

\begin{proof}
The Taylor coefficients of $\Psi$ are nonnegative, and its linear
coefficient is $2/\pi$. Composition preserves nonnegativity and gives
linear coefficient $c_r=(2/\pi)^r$. The Schur product theorem therefore
implies
\begin{equation}\label{gpupper:kernelminorant}
 K_r\succeq c_r\rho,\qquad
 \rho(x,y)=\frac{1+x\cdot y}{d+1}.
\end{equation}
Indeed, the remaining positive Taylor terms are positive semidefinite
entrywise powers of $\rho$. Thus we may realize the Gaussian vector as
\[
 G(x)=\sqrt{\frac{c_r}{d+1}}
       \left(Z_0+\sum_{i=1}^d Z_i x_i\right)+V(x),
\]
where $Z_0,\ldots,Z_d$ are independent standard normals and $V$ is an
independent Gaussian vector with covariance $K_r-c_r\rho$.

Shift only the $Z_i$ so that the output mean becomes
\[
 h(x)=A\left(\sum_{i=1}^d x_i-d+1\right).
\]
The sign of $h$ is $f_\wedge$, and $f_\wedge(x)h(x)\ge A$ everywhere.
The relative entropy of the shifted joint Gaussian law with respect to
the original joint law is
\begin{equation}\label{gpupper:shiftentropy}
 \frac12\frac{d+1}{c_r}
 \left(A^2(d-1)^2+dA^2\right)=\frac12D_{d,L,a}.
\end{equation}
This construction also proves the usual squared reproducing-kernel
norm bound for $h$, without needing to invert the full kernel matrix.

Every centered output has variance one. With
$T=\sqrt{2\log(4m)}$, a Gaussian tail bound and a union bound give
\[
 \mathbb P\{\max_x|G(x)|\ge T\}
 \le2m e^{-T^2/2}=\frac12.
\]
Hence the shifted law gives probability at least $1/2$ to
$\mathcal E=\{f_\wedge(x)G(x)>a\text{ for every }x\}$.
The strict inequalities cause no difficulty: each boundary is a scalar
Gaussian equality event and has probability zero.

Let $p$ and $q\ge1/2$ be the probabilities of $\mathcal E$ under the
original and shifted laws. Relative entropy decreases under the map
that records membership in $\mathcal E$, so
\[
 \frac12D_{d,L,a}
 \ge q\log\frac qp+(1-q)\log\frac{1-q}{1-p}
 \ge q\log\frac1p-\log2
 \ge\frac12\log\frac1p-\log2.
\]
This proves~\eqref{gpupper:rawmass}. Conditioning on the union of all
full-margin function events can only increase the probability of
$\mathcal E$, proving~\eqref{gpupper:targetcost}.
\end{proof}

\begin{corollary}[Every-order online guarantee, including finite width]
\label{gpupper:online}
For the limiting prior, let $B=D_{d,L,a}+\log4$. Against the AND target,
on every nonanticipating input-selection strategy with labels revealed
after each prediction, ordinary consistency Bayes has
\begin{itemize}
\item expected total fresh posterior-sampling mistakes at most $B$;
\item at most $B/\log2$ posterior-majority mistakes;
\item at most $B/\log(3/2)$ rounds on which the correct label has
posterior probability at most $2/3$.
\end{itemize}
These include infinite streams. For the actual Gaussian tanh networks
at $\gamma=0$, the same statements hold, with $B$ replaced by
$D_{d,L,a}+\log8$, for every sufficiently large finite width $N$.
The width threshold can depend on $d,L,a$ but not on the input history,
its order, or its length.
\end{corollary}

\begin{proof}
For any fixed prior $\pi$ and realizable target $f$, let $p_t$ be the
posterior probability of its correct label at the round-$t$ query.
Exact conditioning telescopes to
\[
 \sum_{t=1}^T-\log p_t\le-\log\pi(f).
\]
The sampling bound follows from $1-p_t\le-\log p_t$ and monotone
convergence. A majority mistake has $p_t\le1/2$, while a failure of
correct-label probability strictly greater than $2/3$ contributes at
least $\log(3/2)$.

For fixed $d,L$, the unconditioned finite-width output vector converges
in distribution to $G$ by the hidden-kernel and Gaussian-readout limit
in~\cite{nn-subcritical}. The boundary of the full target margin event
has limiting probability zero. Its raw probability therefore converges
to a value at least $e^{-D_{d,L,a}}/4$, and is at least
$e^{-D_{d,L,a}}/8$ for all sufficiently large $N$. Global margin
conditioning again increases the target probability. Applying the same
fixed-prior potential identity proves all three bounds simultaneously
for every history; no posterior convergence uniform over histories is
being assumed.
\end{proof}

\paragraph{Other fixed Gaussian scales and negative exponents.}
The same argument applies to $G\sim\mathcal N(0,\sigma^2K_r)$ for any
fixed $\sigma>0$. Replace the budget in~\eqref{gpupper:budget} by
\begin{equation}\label{gpupper:scalebudget}
 D_{d,L,a,\sigma}
 =\left(\frac a\sigma+\sqrt{2\log(4m)}\right)^2
   (d+1)(d^2-d+1)\left(\frac\pi2\right)^r.
\end{equation}
Multiplying all layer standard deviations by this fixed $\sigma$
leaves the saturated hidden-kernel limit unchanged and scales the
readout covariance by $\sigma^2$, so the finite-width transfer is valid
for this model as well. For the $\gamma<0$ limit from~\cite{nn-subcritical},
the global margin disappears after output normalization. Set $a=0$ in
the argument; the sign-only limiting prior and all sufficiently large
finite widths obey the corresponding bound. The original finite-width
margin becomes a vanishing normalized threshold, and the limiting
Gaussian assigns zero probability to its boundaries.

These bounds are polynomial in $d$ for each fixed depth and fixed
positive scale, but contain the factor $(\pi/2)^{L-1}$. They therefore
do not rule out an exponential AND lower bound at depth growing with
$d$. Nor is the finite-width threshold quantified here or uniform as
$\sigma\downarrow0$. This is a width-first statement and does not
interchange the strong-prior order of limits.

\section{An exponential conjunction lower bound at growing depth}
\label{sec:gp-deep-and}

The unchanged Gaussian learner can make exponentially many sampled
mistakes on a conjunction when depth grows with dimension. This result
uses the weak-prior scaling $\gamma=0$. It does not address the
fixed-depth strong-prior or $\gamma\ge1$ regimes.

Throughout this section, $d\ge2$, $\mathcal X=\{-1,1\}^d$, and $m=2^d$.
We use tanh hidden layers, one appended constant input, no later biases,
and the independent Gaussian variances
\begin{equation}\label{eq:gp-lower-variances}
 \operatorname{Var}(W^1_{ij})=\frac Nd,\qquad
 \operatorname{Var}(W^s_{ij})=1\quad(1<s<L),\qquad
 \operatorname{Var}(v_i)=\frac1N.
\end{equation}
The prior is conditioned on $|z_w(x)|>1/2$ at every cube input, and
training retains every exact label. Predictions use a fresh posterior
sample. Define
\begin{equation}\label{eq:gp-lower-depth}
 \begin{split}
 k_0(d)&=\left\lceil\frac{\log(d+1)}{\log(4/3)}\right\rceil,\\
 k_1(d)&=\left\lceil\frac{\log(50m^2)}{\log(3/2)}\right\rceil,\qquad
 L_d=1+k_0(d)+k_1(d).
 \end{split}
\end{equation}
In particular, $L_d=O(d)$.

\begin{theorem}[Growing-depth Gaussian NNBayes on the big AND]
\label{thm:gp-deep-and}
Let $f_{\mathrm{AND}}$ be positive only at $(1,\ldots,1)$ and negative
elsewhere. At depth $L_d$ and the variances
\eqref{eq:gp-lower-variances}, there exists $N_0(d)\ge d$ such that,
for every finite $N\ge N_0(d)$, the following statements hold.

After every realizable label history, the correct-label posterior
probability at each previously unseen input belongs to $[2/5,3/5]$.
Consequently, on every ordering of all $m$ distinct inputs,
\begin{equation}\label{eq:gp-lower-mistakes}
 \mathbb E M_m\ge\frac25\,2^d,\qquad
 \mathbb P\{M_m\le2^d/5\}\le e^{-2^d/20}.
\end{equation}
Every query is also a failure to assign the correct label posterior
probability strictly greater than $2/3$. The ordering may be fixed in
advance or selected adaptively from previously unseen inputs, before
the current private posterior draw.

The target has a same-depth tanh implementation of width $d$, with
every coefficient bounded by $d$ and strict signed output margin
greater than $1/2$.
\end{theorem}

The width threshold is only asserted to exist; no quantitative bound
on $N_0(d)$ is proved. The mistake statements concern posterior
sampling and posterior confidence. They do not give a
posterior-majority mistake lower bound.

\begin{lemma}[The fixed-dimensional width limit]
\label{lem:gp-lower-width-limit}
For each fixed $d,L$, under the raw prior
\eqref{eq:gp-lower-variances}, the vector of outputs converges in law to
a centered Gaussian vector with covariance
\begin{equation}\label{eq:gp-lower-kernel}
 K_L(x,y)=\Psi^{\circ(L-1)}
       \left(\frac{1+x\cdot y}{d+1}\right),
 \qquad \Psi(t)=\frac2\pi\arcsin(t).
\end{equation}
\end{lemma}
\begin{proof}
For hidden layer $s$, write
$Q_N^s(x,y)=N^{-1}\sum_{i=1}^N h_i^s(x)h_i^s(y)$.
First-layer rows are independent Gaussian fields with covariance
$N(1+x\cdot y)/d$ before applying tanh. As their scale tends to
infinity, tanh converges almost surely to the sign of a coupled
Gaussian coordinate. The Gaussian sign-correlation identity gives
the limiting second moment
$\Psi((1+x\cdot y)/(d+1))$. Boundedness of the summands gives convergence
of the empirical second moments in probability, simultaneously over
the finitely many input pairs.

Conditionally on a preceding layer, the next preactivation rows are
independent Gaussian fields with covariance $NQ_N^{s-1}$. If the
preceding empirical covariance tends to a fixed matrix with diagonal
one, the corresponding tanh second moments converge to its entrywise
Gaussian sign correlations. This follows by a continuous
Gaussian-square-root coupling and dominated convergence; each
limiting scalar coordinate is nondegenerate. Conditional variances
of the empirical second moments are at most $1/N$.
Induction over the fixed number of layers therefore gives
$Q_N^{L-1}\to K_L$ in probability. Finally, the conditional readout
law is exactly $\mathcal N(0,Q_N^{L-1})$. Convergence of its
characteristic function proves the claim.
\end{proof}

Although the theorem chooses $L=L_d$, this use of
Lemma~\ref{lem:gp-lower-width-limit} is legitimate: dimension and its
chosen finite depth are fixed before width tends to infinity. No
width-limit theorem uniform in increasing depth is required.

\begin{lemma}[Covariance close to the identity]
\label{lem:gp-lower-white}
At the depth \eqref{eq:gp-lower-depth},
\[
 \|K_{L_d}-I_m\|_{\mathrm{op}}\le\eta,\qquad
 \eta=\frac1{100m}.
\]
\end{lemma}
\begin{proof}
For distinct cube vertices,
$|(1+x\cdot y)/(d+1)|\le1-2/(d+1)$.
The map $\Psi$ is odd, increasing, and convex on $[0,1]$.
Since $\Psi(1/2)=1/3$ and $\Psi(1)=1$, its chord on $[1/2,1]$ gives
\[
 1-\Psi(t)\ge\frac43(1-t)\qquad(1/2\le t\le1).
\]
Hence an iterate starting at most $1-2/(d+1)$ must fall below or equal
to $1/2$ within $k_0(d)$ steps: otherwise its gap from one would
be at least $2(4/3)^{k_0(d)}/(d+1)\ge2$, which is impossible.
On $[0,1/2]$, the chord from zero to $(1/2,1/3)$ gives
$\Psi(t)\le(2/3)t$. The next $k_1(d)$ steps thus make every
off-diagonal entry have absolute value at most
\[
 \frac12(2/3)^{k_1(d)}\le\frac1{100m^2}.
\]
All diagonal entries equal one. The absolute row-sum bound for the
operator norm of a symmetric matrix completes the proof.
\end{proof}

\begin{lemma}[Uniform comparison of full-margin orthants]
\label{lem:gp-lower-orthants}
Suppose $G\sim\mathcal N(0,K)$ on $m$ coordinates and
$\|K-I_m\|_{\mathrm{op}}\le1/(100m)$. For each sign vector
$g\in\{-1,1\}^m$, put
\[
 p_g=\mathbb P\{g_iG_i>1/2\text{ for all }i\}.
\]
Then every $p_g$ is positive and
$\max_g p_g/\min_g p_g<4/3$.
\end{lemma}
\begin{proof}
Let $\eta=1/(100m)$, $a=1/2$, and let
$\overline\Phi(t)$ denote the standard Gaussian upper tail.
The eigenvalues of $K$ lie in $[1-\eta,1+\eta]$.
Its density is bounded below by
\[
 (2\pi)^{-m/2}(1+\eta)^{-m/2}
 \exp\left(-\frac{\|z\|^2}{2(1-\eta)}\right)
\]
and above by the same expression with the two occurrences of
$1+\eta$ and $1-\eta$ interchanged.
Integrating over any specified full-margin orthant gives
\[
 \ell\le p_g\le u,
\]
where
\[
 \ell=\left(\frac{1-\eta}{1+\eta}\right)^{m/2}
       \overline\Phi\left(\frac a{\sqrt{1-\eta}}\right)^m,
 \qquad
 u=\left(\frac{1+\eta}{1-\eta}\right)^{m/2}
       \overline\Phi\left(\frac a{\sqrt{1+\eta}}\right)^m.
\]
We bound their ratio without discarding the global-margin condition.
For $\eta\le1/2$,
$\log((1+\eta)/(1-\eta))\le3\eta$.
Both tail arguments lie in $[0,1]$. Writing $\phi$ for the standard
Gaussian density, on that interval
\[
 \frac{\phi(t)}{\overline\Phi(t)}
 \le\frac{\phi(0)}{\int_1^2\phi(s)\,ds}
 \le\frac{\phi(0)}{\phi(2)}=e^2<8.
\]
The difference between the arguments is at most
$a\eta(1-\eta)^{-3/2}\le\sqrt2\,\eta$.
Integrating the logarithmic derivative of $\overline\Phi$ gives
\[
 \log\frac u\ell
 \le (3+8\sqrt2)m\eta<15m\eta=\frac3{20}
 <\log\frac43.
\]
The final inequality follows, for example, from
$\log(1+1/3)\ge1/4$.
\end{proof}

\begin{proof}[Proof of Theorem~\ref{thm:gp-deep-and}]
Apply Lemmas~\ref{lem:gp-lower-white} and
\ref{lem:gp-lower-orthants} to the Gaussian width limit.
For each Boolean function $g$, define its raw neural probability
\[
 p_{N,g}=\mathbb P\{g(x)z_N(x)>1/2\text{ for every }x\}.
\]
The limiting Gaussian is nonsingular, so the boundary of each such
orthant has probability zero. Lemma~\ref{lem:gp-lower-width-limit}
therefore gives $p_{N,g}\to p_g>0$.
There are finitely many Boolean functions at fixed dimension.
The strict limiting ratio bound implies the existence of $N_0(d)$
such that, for every $N\ge N_0(d)$,
\begin{equation}\label{eq:gp-lower-finite-ratio}
 \frac{\max_g p_{N,g}}{\min_g p_{N,g}}\le\frac32.
\end{equation}
Global-margin conditioning divides all these masses by one common
normalizer, preserving their ratios.

After $s$ distinct inputs have been labeled, exactly
$2^{m-s-1}$ full Boolean functions extend the history with either
specified label at a new input. Equation~\eqref{eq:gp-lower-finite-ratio}
therefore places the posterior probability of either label in
$[2/5,3/5]$. This argument is uniform over all histories and does
not depend on their probability or on the target. Conditional
expectation gives the first bound in \eqref{eq:gp-lower-mistakes}.

For the second bound, let $I_t$ be the sampled mistake indicator.
Conditionally on the past and the current previously unseen input,
\[
 \mathbb E[2^{-I_t}\mid\text{past, current input}]
 =1-\frac12\mathbb P(I_t=1\mid\text{past, current input})
 \le\frac45.
\]
Iterating and applying Markov's inequality gives
\[
 \mathbb P\{M_m\le m/5\}
 \le 2^{m/5}(4/5)^m\le e^{-m/20},
\]
using $\log2<3/4$ and $\log(4/5)\le-1/5$.
The confidence claim follows from $3/5<2/3$.

Finally, realize $f_{\mathrm{AND}}$ with one active neuron in each
hidden layer, padding to width $d$. The first active neuron computes
$\tanh(\sum_jx_j-(d-1))$; each later active neuron computes
$\tanh(2h)$ from the previous one; the readout has coefficient one.
The first signed activation is at least $\tanh(1)>1/2$.
Since $2\tanh(1)>1$, subsequent layers preserve a signed activation
of magnitude at least $\tanh(1)$. Every coefficient has magnitude at
most $d$ for $d\ge2$, and the desired strict margin follows.
\end{proof}

The theorem applies, in particular, to a fixed lexicographic ordering
of the cube. It does not use a circuit tie rule or inspect a sampled
hypothesis before choosing its query. The conjunction has $d-1$
Boolean AND gates and only $O(d+L_d)$ nonzero coefficients in the
displayed neural implementation. Its exponential mistake count is
compatible with fixed-depth polynomial theorems because $L_d$ grows
with $d$, and because \eqref{eq:gp-lower-variances} lies at
$\gamma=0$.

% Width-one geometry and the separately specified anisotropic prior.
% Requires amsmath, amssymb, amsthm; lemma/proposition/corollary environments.
\section{Exact width-one geometry at the start of an AND layer}
\label{sec:and-width-one}

This section concerns a width-one network with one hidden tanh unit,
\[
 z_{a,b}(x)=b\tanh\!\left(a_0+\sum_{i=1}^d a_i x_i\right),
 \qquad x\in\{-1,1\}^d,
 \qquad
 \kappa(a,b)=d\sum_{i=0}^d a_i^2+b^2.
\]
The prior is the centered Gaussian with density proportional to
$\exp[-\kappa(a,b)/(2\sigma^2)]$, conditioned before training on
$|z_{a,b}(x)|>1/2$ at every cube input.
For a Boolean threshold function $g$, define its affine unit-margin cost
\[
 Q(g)=\min\left\{\sum_{i=0}^d a_i^2:
       g(x)\left(a_0+\sum_{i=1}^d a_i x_i\right)\ge1
       \text{ for every }x\right\}.
\]
The minimum exists: the feasible set is a nonempty closed polyhedron
and the objective is coercive. Non-threshold functions are not represented
by this width-one architecture.

\begin{lemma}[Reduction to affine margin cost]\label{lem:and-affine-reduction}
For every represented Boolean function $g$, its minimum strict-margin
neural energy is
\[
 E(g)=\Phi_d(Q(g)),\qquad
 \Phi_d(Q)=\min_{u>0}\left\{dQu^2+\frac14\coth^2u\right\}.
\]
The function $\Phi_d$ is strictly increasing on $Q>0$.
\end{lemma}
\begin{proof}
Replacing $(a,b)$ by $(-a,-b)$ preserves the real output and energy, so
we can take $b>0$. A strict output margin requires $b>1/2$. Put
$u=\operatorname{arctanh}(1/(2b))>0$. The constraints imply
$g(x)(a_0+\sum_i a_i x_i)>u$, so
$\sum_i a_i^2\ge u^2Q(g)$ and $b^2=\frac14\coth^2u$.
Conversely, multiply a unit-margin minimizer for $g$ by $(1+\epsilon)u$
and take $b=\frac12\coth u$. This gives strict output margin, with
energy tending to the displayed expression as $\epsilon\downarrow0$.
This proves equality of the infima. For $Q>0$ the objective tends to
infinity at both endpoints $u\downarrow0$ and $u\to\infty$, so its
minimum is attained at some $u_Q>0$. If $Q_2>Q_1$, then
\[
 \Phi_d(Q_2)
 =dQ_1u_{Q_2}^2+\tfrac14\coth^2u_{Q_2}
   +d(Q_2-Q_1)u_{Q_2}^2
 >\Phi_d(Q_1).
\]
\end{proof}

\paragraph{Small-variance function probabilities.}
For this finite-dimensional Gaussian law, the function event
$U_g=\{(a,b):g(x)z_{a,b}(x)>1/2\ \forall x\}$ satisfies
\begin{equation}\label{eq:affine-rate}
 \lim_{\sigma\downarrow0}\sigma^2\log\mathbb P_\sigma(U_g)
 =-\frac12 E(g).
\end{equation}
Indeed, the lower Gaussian large-deviation bound applies to the open
set $U_g$, and the upper bound applies to its closure. Their energy
infima coincide: any parameter vector with all signed outputs at
least $1/2$ can be approximated from within $U_g$ by multiplying the
readout by $1+\epsilon$. The standard finite-dimensional Gaussian
rate function is half the quadratic energy. The argument holds
unchanged for every positive definite diagonal quadratic energy.
There are only finitely many represented Boolean functions. Therefore,
after conditioning on any fixed label history, all strong-prior
posterior mass concentrates on its minimum-energy functions; if the
minimum is unique, its posterior probability tends to one.

Write $|x|_+=|\{i:x_i=+1\}|$. Fix $2\le k\le d$ and consider the
history that labels $1^d$ positive and every input with $|x|_+<k$
negative. These labels agree with the AND target, which is positive
only at $1^d$. No labels at other inputs are imposed yet.

\begin{lemma}[Essential affine coefficients]\label{lem:and-essential}
In a minimum-squared-norm unit-margin representation of a threshold
function, every nonzero input coefficient is essential to that function,
and every essential input coefficient has absolute value at least one.
\end{lemma}
\begin{proof}
If the function is independent of $x_i$, then, for each fixed assignment
of the other coordinates, both signs of $x_i$ have the same label.
Averaging their two signed margin inequalities shows that setting
$a_i=0$ preserves every unit-margin constraint. It strictly decreases
squared norm if $a_i\ne0$. If $x_i$ is essential, two neighboring cube
points differing only in $x_i$ have opposite labels. Their affine scores
have opposite signs and magnitudes at least one; their difference has
magnitude $2|a_i|$. Hence $|a_i|\ge1$.
\end{proof}

\begin{proposition}[All minimizers at a complete lower-layer prefix]
\label{prop:and-phase-minimizers}
Among all width-one Boolean functions consistent with this history,
the smallest affine unit-margin cost is
\[
 Q_{d,k}=
 \begin{cases}
  2k-1,&d\ge2k-1,\\
  d+(2k-1-d)^2,&d<2k-1.
 \end{cases}
\]
If $d\ge2k-1$, the minimizing functions are precisely
\[
 h_T(x)=+1
 \quad\Longleftrightarrow\quad
 |\{i\in T:x_i=+1\}|\ge k,
 \qquad |T|\in\{2k-1,2k-2\}.
\]
Their minimizing affine scores have coefficient one on $T$, zero
elsewhere, and bias zero for $|T|=2k-1$ or bias $-1$ for $|T|=2k-2$.
If $d<2k-1$, there is exactly one minimizing function: the symmetric
threshold $|x|_+\ge k$, with coefficients one and bias $d-2k+1$.
The corresponding minimum neural energy is $\Phi_d(Q_{d,k})$.
\end{proposition}
\begin{proof}
Choose a minimizing function and its minimum-norm unit-margin affine
representation. Lemma~\ref{lem:and-essential} applies. Let $p$ be the
number of positive input coefficients, write their sum as $P$, and let
$B$ be the sum of absolute values of the negative input coefficients.
In particular $P\ge p$.

One has $p\ge k$: the input which is positive exactly on the positive
coefficient coordinates has affine score at least that at $1^d$, hence
at least one. If $p<k$, the history requires this input to be negative,
a contradiction.

If $p\ge2k-1$, the positive coefficients alone contribute squared norm
at least $p\ge2k-1$. If $p\le2k-2$, average the negative margin
constraints over inputs that are positive on exactly $k-1$ of those
$p$ coordinates and negative everywhere else. This gives
\[
 a_0-P+B+\frac{2(k-1)}pP\le-1.
\]
Since $1-2(k-1)/p\le0$ and $P\ge p$,
\[
 a_0\le\left(1-\frac{2(k-1)}p\right)P-B-1
       \le p-2k+1-B\le p-2k+1<0.
\]
Consequently the total squared norm is at least
\[
 p+(2k-1-p)^2.
\tag{*}
\]
For integer $p\le2k-2$, this is at least $2k-1$, with equality only
at $p=2k-2$. Together with the preceding branch, equality at $2k-1$
forces either $p=2k-1$ with zero bias or $p=2k-2$ with bias $-1$.
It also forces every positive coefficient to equal one and excludes
all negative coefficients. The displayed functions realize these
bounds: all their affine scores are odd integers, and their positive
inputs contain at least $k$ positive coordinates.

If $d<2k-1$, every admissible $p$ lies between $k$ and $d\le2k-2$.
The expression in (*) strictly decreases as integer $p$ increases
in this range. It is therefore minimized only at $p=d$, giving
$d+(2k-1-d)^2$. Equality forces all coefficients to equal one and
the bias to equal $d-2k+1$. This score also takes odd integer values
and realizes the required symmetric threshold. The energy conclusion
follows from Lemma~\ref{lem:and-affine-reduction}.
\end{proof}

\paragraph{Why the circuit conjunction characterization fails here.}
A conjunction of $k$ positive literals has
\[
 Q(\operatorname{AND}_k)=k+(k-1)^2=k^2-k+1.
\]
To verify the lower bound, remove irrelevant coefficients, use
essentiality to make each remaining positive coefficient at least one,
and average the $k$ negative constraints obtained by flipping one of
the $k$ selected coordinates. This forces the bias to be at most
$-(k-1)$. Unit coefficients and that bias attain equality.
For $k\ge3$ and $d\ge2k-1$,
\[
 Q(\operatorname{AND}_k)-(2k-1)=(k-1)(k-2)>0.
\]
Thus every $k$-conjunction is strictly more expensive than the actual
minimizers at this prefix. The globally smallest circuit functions
at the analogous circuit prefix are $k$-conjunctions; the Gaussian
neural energy gives a different answer already in this one-unit model.

\paragraph{Exact masses in this particular minimizing class.}
Equality of minimum energy alone does not imply equal Gaussian function
probabilities. Here, however, an additional exact symmetry does imply it.
Ordinary input permutations make all functions with a fixed support
size equally probable. To connect the two support sizes, take
$j\in T$ with $|T|=2k-1$ and transform the affine parameters by
\[
 a'_0=-a_j,\qquad a'_j=a_0,\qquad a'_i=a_i\quad(i\ne j).
\]
This signed column swap preserves the Gaussian law and energy.
For $x\in\{-1,1\}^d$, set
\[
 y_j=-x_j,\qquad y_i=x_ix_j\quad(i\ne j).
\]
Then $a'_0+\sum_i a'_ix_i=x_j(a_0+\sum_i a_iy_i)$, so oddness of
tanh gives $z'(x)=x_jz(y)$. The map $x\mapsto y$ is a bijection,
and hence the global-margin event is preserved. Applied to majority
on $T$, the transformed sign is
\[
 x_j\operatorname{sign}\!\left(\sum_{i\in T}y_i\right)
 =\operatorname{sign}\!\left(-1+\sum_{i\in T\setminus\{j\}}x_i\right),
\]
which is precisely the threshold of $k$ positives on the even support
$T\setminus\{j\}$. Thus all functions in the minimizing class have
exactly equal prior mass at every positive variance. Equal mass persists
among whichever of them survive later label conditioning.

By the fixed-width Gaussian rate theorem and finiteness of the Boolean
function class, the strong-prior posterior at this prefix converges to
uniform sampling from these minimizers. At positive variance it also
assigns mass to higher-energy functions; it is not exactly this uniform
minimum-class rule. In other histories, equal energies need not be
accompanied by the symmetry just proved.

\begin{corollary}[The first minimum-energy tier cannot itself cause an
exponential phase]\label{cor:and-first-tier}
Suppose $d\ge2k-1$ and start from the complete lower-layer prefix above.
Let
\[
 M_{d,k}=\binom d{2k-1}+\binom d{2k-2}.
\]
For the width-one strong-prior limiting posterior, the expected number
of sampled mistakes before this minimum-energy tier is exhausted is at
most $1+\log M_{d,k}\le1+(d+1)\log2$. The bound includes the round which exhausts
it. It holds under every adaptive nonanticipating continuation. If
$d<2k-1$ and $k<d$, its sole initial minimum is eliminated by the first
weight-$k$ negative example, causing one mistake on that example. If
$k=d$, the prefix has already revealed the entire cube and no further
mistakes are possible.
\end{corollary}
\begin{proof}
While at least one of the original minimizers survives, it still has
energy $\Phi_d(Q_{d,k})$, so every strong-prior posterior is uniform
on the surviving original minimizers. Eliminating a fraction $q<1$
of a uniform finite class reduces its logarithmic size by at least $q$.
The final exhaustion costs at most one. Summing these conditional error
probabilities and taking expectations proves the claim. The singleton
case follows from the unique symmetric threshold's positive prediction
on every weight-$k$ input.
\end{proof}

\paragraph{Limits of the conclusion.}
The proposition characterizes the first energy tier at a complete
lower-layer prefix. After that tier is exhausted, higher-energy
functions can become optimal; their complete sequence is not
characterized here. Thus the corollary is not an upper bound for the
whole Hamming layer or for the entire AND-learning process.
All statements concern width one, one hidden tanh unit, and global
margin conditioning. For a growing AND target, width one need not
satisfy the separate theorem convention $N\ge C_L(f)$, because that
complexity also bounds the magnitude of every teacher weight. No
extension to all admissible widths, and no exponential lower bound
for the original Gaussian neural posterior, is asserted by this width-one calculation.

\subsection{An arbitrarily small anisotropy changes the width-one answer}
\label{sec:and-anisotropy}

The preceding equality of masses uses exact symmetry. The next result
changes the prior and is not a lower bound for the canonical isotropic
first-layer law. It also retains width one; it does not satisfy the
separate $N\ge C_L(f)$ convention for the growing AND target in general.
It does show that a small independent-Gaussian perturbation can turn the
same concrete neural minimizing class into an exponential mistake path.

\begin{proposition}[Near-isotropic independent-Gaussian variant]
\label{prop:and-anisotropic}
Let $k\ge2$ and $d\ge8k$. For every $\eta>0$, there are coefficients
$1\le\lambda_i\le1+\eta$, $0\le i\le d$, with the following property.
Use a width-one prior with independent centered Gaussian weights
\[
 \operatorname{Var}(a_i)=\frac{\sigma^2}{d\lambda_i},\qquad
 \operatorname{Var}(b)=\sigma^2,
\]
and the same all-input output-margin conditioning. Its quadratic energy is
\[
 \kappa_\lambda(a,b)=d\sum_{i=0}^d\lambda_i a_i^2+b^2.
\]
For the AND target, there is a fixed distinct-input order on
which the strong-prior posterior makes at least
\[
 \frac{\binom d k}{\binom{2k-1}k}\ge2^k
\]
posterior-sampling mistakes with probability one. Posterior majority is
wrong on the same rounds.

For every fixed $0<\zeta<1/2$, one sufficiently small positive variance
makes the actual Gaussian posterior's wrong-label probability at least
$1-\zeta$ on every one of these rounds. Its expected sampled mistakes
on them are therefore at least $(1-\zeta)2^k$, and its majority is
wrong on every one. In particular, $k=\lfloor d/8\rfloor$ gives
an exponential lower bound for this perturbed width-one model.
\end{proposition}
\begin{proof}
Start with the history of Proposition~\ref{prop:and-phase-minimizers}.
Let $\mathcal C$ be its minimizing class and $Q_0=2k-1$. Every other
consistent threshold function has strictly larger isotropic affine cost.
There are finitely many threshold functions, so their costs exceed $Q_0$
by a positive common gap $\Delta$.

For $\epsilon>0$, put $\lambda_i=1+\epsilon v_i$, where the numbers
$v_i\in(0,1]$ have distinct sums over distinct subsets of
$\{0,\ldots,d\}$; for example $v_i=2^{-i}$ has this property.
For each fixed threshold function $h$, let
\[
 Q_h(\epsilon)=\min\left\{\sum_i(1+\epsilon v_i)a_i^2:
             h(x)(a_0+\textstyle\sum_{i\ge1}a_ix_i)\ge1
             \text{ for all }x\right\}.
\]
Its isotropic optimizer $a_h$ is unique by strict convexity. The
optimizers at $\epsilon\downarrow0$ converge to $a_h$: boundedness
follows from comparison with $a_h$, and every limit point is an isotropic
optimizer. Comparing the perturbed and unperturbed objectives in both
directions consequently gives
\[
 Q_h(\epsilon)=Q_h(0)+\epsilon\sum_i v_i(a_h)_i^2+o(\epsilon).
\]
For $h\in\mathcal C$, the nonzero squared coordinates of $a_h$ are
exactly the indicator of its odd support $T$, or of its even support
$T$ together with the bias index $0$. These augmented supports are
pairwise different. Their derivative sums are therefore pairwise
distinct. Since $\mathcal C$ is finite, all its perturbed costs are
pairwise distinct for sufficiently small positive $\epsilon$.

For every $h\in\mathcal C$,
$Q_h(\epsilon)\le(1+\epsilon)Q_0$. All other initially consistent
functions have $Q_h(\epsilon)\ge Q_0+\Delta$, because $v_i>0$.
Taking also $\epsilon<\Delta/Q_0$ and $\epsilon\le\eta$ ensures that
every member of $\mathcal C$ remains cheaper than every outsider.
The affine-to-tanh reduction is unchanged, with $Q_h(\epsilon)$ in
place of $Q(h)$, so the same ordering holds for neural energies.

Until $\mathcal C$ is exhausted, the currently consistent function of
minimum energy is therefore a unique member of $\mathcal C$. Choose
any $k$-element subset of its support and query the input positive
exactly on that subset. The selected function predicts positive; the
AND target is negative because $k<d$. The query is new, since a
consistent selected function cannot be positive on a previously
revealed negative query. This step eliminates exactly those candidate
supports containing the queried $k$-set.

There are $\binom d{2k-1}$ candidates with odd support size $2k-1$.
A single $k$-set query eliminates at most
$\binom{d-k}{k-1}$ of these. Exhausting $\mathcal C$ therefore
requires at least
\[
 \frac{\binom d{2k-1}}{\binom{d-k}{k-1}}
 =\frac{\binom d k}{\binom{2k-1}k}
 \ge\frac{(d/k)^k}{4^k}\ge2^k
\]
queries, each a certain mistake in the strong-prior limit.
Use a fixed convention for choosing each $k$-set; the complete query
sequence can then be computed before the run and completed to a
permutation of the cube. It does not inspect any current random draw.

Finally, at each of these finitely many histories the minimum-energy
function is unique with a positive energy gap. The Gaussian rate
argument applies to every positive definite diagonal quadratic energy.
Thus, as the common variance tends to zero, the posterior mass of that
function tends to one. Taking the minimum of finitely many positive
variance thresholds makes its wrong-label probability at least
$1-\zeta$ at every selected round simultaneously. The input sequence
and labels are fixed, so posterior randomness at earlier rounds does
not alter this argument.
\end{proof}

\begin{corollary}[An admissible width-one majority target]
\label{cor:majority-anisotropic}
Let $d\ge17$ be odd and put $k=\lfloor d/8\rfloor$. The conclusion
of Proposition~\ref{prop:and-anisotropic} also holds for the target
\[
 f_{\rm maj}(x)=\operatorname{sign}\!\left(\sum_{i=1}^d x_i\right).
\]
This target has bounded-weight teacher complexity $C_2(f_{\rm maj})=1$,
so the learner width $N=1$ satisfies $N\ge C_2(f_{\rm maj})$.
For every $\eta>0$, an independent-Gaussian first-layer covariance
perturbation with $1\le\lambda_i\le1+\eta$ therefore admits a fixed
input permutation with at least $2^{\lfloor d/8\rfloor}$ certain
strong-prior mistakes. At all sufficiently small positive variances,
the expected sampled mistakes on those rounds are at least
$(1-\zeta)2^{\lfloor d/8\rfloor}$ for any prescribed
$0<\zeta<1/2$, and posterior majority is wrong on all of them.
\end{corollary}
\begin{proof}
The network $z(x)=\tanh(\sum_i x_i)$ has one hidden unit, zero bias,
and all input and readout weights equal to one. Since $d$ is odd,
$|\sum_i x_i|\ge1$ on the cube, so its signed output margin is at
least $\tanh(1)>1/2$. Thus its bounded-weight teacher complexity is
one.

The construction in Proposition~\ref{prop:and-anisotropic} uses only
the positive label at $1^d$ and negative labels at inputs of Hamming
weight at most $k$. Because $k<d/2$, all these labels agree with
$f_{\rm maj}$. Consequently every posterior and every forced query
in that construction is unchanged when the target is $f_{\rm maj}$.
After that finite sequence, reveal any remaining inputs with their
majority labels. The same Gaussian prior, energy gaps, and positive
variance thresholds prove the stated bounds.
\end{proof}

\paragraph{Scope of the perturbed examples.}
The parameters of the width-one target
$\tanh(\sum_i x_i-d+1)$ have magnitude $O(d)$ and strict output margin
$\tanh(1)>1/2$. Thus the example has a simple, polynomially bounded
neural teacher, but this AND realization does not establish the
bounded-weight width-admissibility convention.
Corollary~\ref{cor:majority-anisotropic} removes that width issue for
a different simple target. Both examples change the exact coordinate
symmetry of the first layer. Neither is an exponential lower bound
for the canonical isotropic prior; no extension of the AND lower
bound to all admissible larger learner widths is established.

\paragraph{High probability and ordinary iid learning.}
The finite-variance lower bound can also force all the designated
mistakes with probability at least $1-\delta$: if there are $T$ forced
rounds, apply the proposition with $\zeta<\delta/T$ and use a union
bound. Choosing also $\zeta<1/3$ makes every designated round a failure
of the strict $2/3$ confidence requirement.

The same model retains polynomial iid sample complexity. Its supported
Boolean functions are affine threshold functions. The arrangement of
$m=2^d$ homogeneous hyperplanes in the $(d+1)$-dimensional affine
coefficient space gives at most
$2\sum_{j=0}^d\binom{m-1}{j}\le2(d+1)m^d$
possible strict sign patterns. For iid examples from any fixed input
distribution, a union bound therefore implies that every consistent
function has error at most $\varepsilon$, except with probability
$\delta$, once
\[
 n\ge\frac{\log\bigl(2(d+1)2^{d^2}\bigr)+\log(1/\delta)}
              {\varepsilon}.
\]
This controls posterior-sampling test error for every prior on this
class. The test-distribution probability of inputs failing strict $2/3$ confidence is at
most three times the posterior-averaged error. Thus the anisotropic
example separates polynomial iid learning from exponential online
mistakes under an adversarially chosen fixed order; it does not merely
use a model that was already statistically hard under iid data.


\section{What the results do and do not settle}
The depth-$O(d)$ lower bound is a counterexample for the actual
independent-Gaussian neural prior. The inputs need no adversarial
ordering. Its proof controls full-margin orthant probabilities, not
merely unconditioned covariance or Gaussian total variation, so the
initial support restriction is retained. It concerns a weak prior;
it is not a counterexample to the fixed-depth circuit-prior guarantees
for sufficiently small variance, above mean field, or mean field.

At fixed depth in the Gaussian-process regime, the explicit AND
upper bound rules out an exponential asymptotic lower bound once width
is sufficiently large. It pays a factor $(\pi/2)^{L-1}$ and depends on
a fixed common Gaussian scale when such a scale is included. The bound
is not uniform in a scale tending to zero, and the eventual-width
statement is not an all-width quantitative theorem.

For the equal-coordinate-variance strong-prior learner, and for the
$\gamma\ge1$ family, an exponential arbitrary-order AND lower bound
remains unproved. The neural phase calculation identifies a failure
of the circuit proof, not a positive theorem for all subsequent phases
or all widths. Conversely, an elementary upper bound through
$-\log\pi_N(f_\wedge)$ becoming large does not prove many mistakes.
The anisotropic one-neuron example shows that an arbitrarily small
change in coordinate variances can give exponential mistakes even for
a unit-weight, width-one majority teacher, while retaining polynomial
iid sample complexity. The variance perturbation is fixed before
the common prior scale is made sufficiently small. This does not
justify interchanging those two limits or asserting the result at
exactly equal coordinate variances.

The structured-input guarantees of~\cite{neural-online} remain valid
and do not cover arbitrary input sequences.

\paragraph{Verification.}
The new results have conventional proofs and independent AI-agent
mathematical reviews. They are not Lean formalized. Earlier Lean
coverage of the Gaussian width limits does not automatically verify
these new consequences. The accompanying verification record specifies
the checked statements, source hashes, and PDF build checks.

\begin{thebibliography}{9}
\raggedright
\bibitem{nn-subcritical}
ChatGPT-6 Astra, Kaarel H\"anni, and Dmitry Vaintrob.
\emph{Exponential parity sample complexity below the mean-field scaling}.
\url{https://github.com/kaarelh/AI-works/tree/main/neural-network-generalization/main-results/subcritical-exponential}.
\bibitem{occam}
ChatGPT-6 Astra, Claude, Kaarel H\"anni, and Dmitry Vaintrob.
\emph{Exponential mistakes and polynomial guarantees for circuit Occam learning}.
\url{https://github.com/kaarelh/AI-works/tree/main/neural-network-generalization/works-in-progress/online-mistakes/occam-distribution-shift}.
\bibitem{neural-online}
ChatGPT-6 Astra, Claude, Kaarel H\"anni, and Dmitry Vaintrob.
\emph{Online mistake bounds for Gaussian neural networks and priors preserving their function rankings}.
\url{https://github.com/kaarelh/AI-works/tree/main/neural-network-generalization/works-in-progress/online-mistakes/neural-function-priors}.
\end{thebibliography}
\end{document}
